This appendix proves every result in the paper. Statements use the class model of Section~\ref{sec:model}. A
statement that needs the example says so. Each equilibrium proof states the equilibrium tuple of
Definition~\ref{def:equilibrium}, lists its conditions, and shows that the fixed tuple satisfies each of them.

% STRUCTURE: Definition 1 and the conditions (E1)-(E6) <- Section 2.4; eta <- Section 2.3
\paragraph{The equilibrium tuple.}
Write \(\eta(-1)=\eta(1)=\beta\) and \(\eta(0)=1-2\beta\) for the liquidity order distribution of
Section~\ref{sec:slot}.

\begin{definition}[Equilibrium]\label{def:equilibrium}
An \emph{equilibrium} is a tuple
\[
\mathcal{E}=\bigl(\sigma_H,\sigma_L,\ \mu,\ P,\ \bar\mu,\ \zeta\bigr),
\]
together with the bids of the prepared bidders. Its components are as follows.
\begin{itemize}[leftmargin=2em]
\item \(\sigma_j=\bigl(\sigma_j(-1),\sigma_j(0),\sigma_j(1)\bigr)\) is a probability distribution on
\(\{-1,0,1\}\), the order of the informed trader in class \(j\in\{H,L\}\). Mixed and wrong-signed orders are
allowed.
\item \(\mu:\{-1,0,1\}\to[0,1]\) is the market makers' posterior \(\mu(x)=\Pr(j=H\mid X=x)\).
\item \(P:\{-1,0,1\}\to\mathbb{R}\) is the price after each order.
\item \(\bar\mu\) assigns to each price \(\bar P\) in the image of \(P\) the rival's belief
\(\Pr(j=H\mid P(X)=\bar P)\).
\item \(\zeta(\bar P)\in[0,1]\) is the probability that a rival with the high cost prepares after the price
\(\bar P\). A rival with the zero cost prepares after every price.
\end{itemize}
For each order \(x\), write \(f_j(x)=(1-\lambda)\eta(x)+\lambda\sigma_j(x)\) for its probability in class
\(j\), and \(e(x)=\rho+(1-\rho)\,\zeta(P(x))\) for entry after it. The tuple satisfies six conditions.
\begin{enumerate}[label=(E\arabic*),leftmargin=3em]
\item \emph{Bayes consistency.} For every order \(x\),
\(\mu(x)=\pi f_H(x)/[\pi f_H(x)+(1-\pi)f_L(x)]\).
\item \emph{Beliefs of the rival.} For every price \(\bar P\) in the image of \(P\), \(\bar\mu(\bar P)\) is the
posterior of class \(H\) given the event \(\{x:P(x)=\bar P\}\), computed from \(f_H\) and \(f_L\).
\item \emph{Entry optimality.} For every price \(\bar P\) in the image, \(\zeta(\bar P)=1\) if
\(B(\bar\mu(\bar P))>c_H\) and \(\zeta(\bar P)=0\) if \(B(\bar\mu(\bar P))<c_H\). At
\(B(\bar\mu(\bar P))=c_H\), the rival's tie rule applies.
\item \emph{Competitive prices.} For every order \(x\), \(P(x)=\E[V_T\mid X=x]\).
\item \emph{Optimality of the informed trader.} For each class \(j\), \(\sigma_j(q)>0\) only if \(q\) maximizes
\[
U_j(q)=q\,\bigl[t_0+e(q)(t_j-t_0)-P(q)\bigr]-k|q|
\]
over \(q\in\{-1,0,1\}\), with \(P(\cdot)\) and \(e(\cdot)\) held at their values in \(\mathcal{E}\).
\item \emph{Truthful bids.} Each prepared bidder whose value meets the reserve bids its value.
\end{enumerate}
\end{definition}

The gross profit of a zero-cost rival is an expected auction rent, which is never negative. Preparation is
therefore optimal for it after every price, and the definition fixes this choice. The probability \(f_j(x)\) is
at least \((1-\lambda)\eta(x)>0\) under every strategy. Every order therefore occurs with positive probability in
both classes. Bayes' rule defines \(\mu\) after every order, and no off-path belief enters the definition.

The two indifferences in (E3) and (E5) follow the default tie rules of Section~\ref{sec:equilibrium}: an
indifferent informed trader may mix over its best orders, and an indifferent high-cost rival prepares.
Appendix~\ref{app:coexistence} states the alternatives, under which an indifferent informed trader abstains or an indifferent
high-cost rival may prepare with any probability.

\subsection{Acquisition payoffs}\label{app:payoffs}

\paragraph{Truthful bids.}
Fix a bidder with value \(\theta\), and let \(Y\ge0\) be the highest competing admissible bid, with \(Y=0\) when
there is none. Any admissible bid that wins pays \(\max\{p,Y\}\), so a bid yields either \(\theta-\max\{p,Y\}\) or
zero. Suppose \(\theta\ge p\), and consider the bid \(\theta\). If \(\theta>\max\{p,Y\}\), it wins and earns
\(\theta-\max\{p,Y\}>0\). If \(\theta<\max\{p,Y\}\), then \(Y>\theta\), and the bid loses and earns zero. If
\(\theta=\max\{p,Y\}\), a win pays \(\theta\), so the bid earns zero.

The bid \(\theta\) therefore earns \((\theta-\max\{p,Y\})_+\), the larger of the two possible payoffs, at every \(Y\).
Suppose \(\theta<p\). Every bid that wins pays at least \(p>\theta\), so no bid earns more than zero, the payoff of no
bid. Truthful bids, and no bid below the reserve, are therefore weakly dominant at every competing bid. Hence
condition (E6) of Definition~\ref{def:equilibrium} holds in every tuple. No payoff below rests on a belief.

\paragraph{Proof of Proposition~\ref{prop:returns}.}
Without entry the target receives \(p\1\{R\ge p\}\). A prepared rival with \(\theta\ge p\) wins at
\(\theta>R\) and pays \(\max\{p,R\}\). A loss makes the incumbent pay \(\theta\). The target receives
\(\max\{p,\min(R,\theta)\}\), and the rival keeps \((\theta-\max\{p,R\})_+\). At \(\theta<p\) the
rival does not bid. The target again receives \(p\1\{R\ge p\}\), and the rival keeps nothing.

\emph{The example.} Let \(R\) be uniform on \([0,r]\) with \(0<p<\ell<r<h\), as in \eqref{eq:order}. Then
\(t_0=\Pr(R\ge p)\,p=p(1-p/r)\). Since \(R\le r<h\), the high-class rival always wins and pays
\(\max\{p,R\}\). The low-class rival wins and pays \(\max\{p,R\}\) when \(R<\ell\); it loses when
\(R\ge\ell\). Integration over \(R\) gives
\begin{equation}\label{eq:a-payoffs}
\begin{aligned}
t_H&=\frac1r\Bigl[\int_0^p p\,dy+\int_p^r y\,dy\Bigr]
=\frac1r\Bigl[p^2+\frac{r^2-p^2}2\Bigr]=\frac r2+\frac{p^2}{2r},\\
t_L&=\frac1r\Bigl[\int_0^p p\,dy+\int_p^\ell y\,dy+\int_\ell^r\ell\,dy\Bigr]
=\frac1r\Bigl[p^2+\frac{\ell^2-p^2}2+\ell(r-\ell)\Bigr]=\ell-\frac{\ell^2-p^2}{2r},\\
g_H&=\frac1r\int_0^r\bigl(h-\max\{p,y\}\bigr)\,dy=h-t_H=h-\frac r2-\frac{p^2}{2r},\\
g_L&=\frac1r\Bigl[\int_0^p(\ell-p)\,dy+\int_p^\ell(\ell-y)\,dy\Bigr]
=\frac1r\Bigl[p(\ell-p)+\frac{(\ell-p)^2}2\Bigr]=\frac{\ell^2-p^2}{2r}.
\end{aligned}
\end{equation}

The spread is
\[
\DT(r)=t_H-t_L=\frac r2+\frac{p^2}{2r}-\ell+\frac{\ell^2-p^2}{2r}=\frac{r^2-2\ell r+\ell^2}{2r}
=\frac{(r-\ell)^2}{2r}.
\]
The sums are \(t_H+g_H=h\) and \(t_L+g_L=\ell\), as the split requires. These are the expressions in \eqref{eq:payoffs}.

Differentiation in \(r\) gives
\[
\begin{gathered}
\DT'(r)=\frac{2(r-\ell)\cdot2r-2(r-\ell)^2}{4r^2}=\frac{(r-\ell)(r+\ell)}{2r^2}=\frac12-\frac{\ell^2}{2r^2},\\
g_H'(r)=-\frac12+\frac{p^2}{2r^2},\qquad g_L'(r)=-\frac{\ell^2-p^2}{2r^2}.
\end{gathered}
\]
The signs in \eqref{eq:derivatives} follow from \(r>\ell\), \(r>p\), and \(\ell>p\). For every
\(\mu\in[0,1]\), \(\partial B_r(\mu)/\partial r=(1-\mu)g_L'(r)+\mu g_H'(r)<0\).

\emph{General classes and a continuous \(F\).} Write the payments as integrals of indicators. For
\(\theta\ge p\),
\begin{equation}\label{eq:a-indicators}
\max\{p,\min(R,\theta)\}=p+\int_p^\theta\1\{R>y\}\,dy,
\qquad
(\theta-\max\{p,R\})_+=\int_p^\theta\1\{R<y\}\,dy.
\end{equation}
The two indicators sum to one except at the single point \(y=R\), which gives \(t(\theta)+g(\theta)=\theta\).

Every integrand in \eqref{eq:a-indicators} is nonnegative and jointly measurable in \((R,y)\). Tonelli's theorem
therefore moves the expectation over \(R\) inside each integral. Since \(F\) is continuous,
\(\Pr(R>y)=1-F(y)\). This gives \(t(\theta)=p+\int_p^\theta[1-F(y)]\,dy\) and \(g(\theta)=\int_p^\theta
F(y)\,dy\) for \(\theta\ge p\). For \(\theta<p\), \(t(\theta)=t_0=p[1-F(p)]\) and \(g(\theta)=0\).

At \(\theta=p\), \(t\) jumps by \(pF(p)\). When the incumbent fails the reserve, an admissible rival buys at
\(p\), and an inadmissible one leaves the target unsold. Hence, for every \(\theta\),
\[
t(\theta)=t_0+pF(p)\1\{\theta\ge p\}+\int_p^\infty\1\{y\le\theta\}[1-F(y)]\,dy,\qquad
g(\theta)=\int_p^\infty\1\{y\le\theta\}F(y)\,dy.
\]

Take the class expectation and apply Tonelli's theorem again, with \(\E_j\1\{\theta\ge p\}=1-G_j(p^-)\) and
\(\Pr_j(\theta\ge y)=1-G_j(y^-)\). The functions \(G_j(y^-)\) and \(G_j(y)\) differ at countably many \(y\), a
Lebesgue-null set, so the integral is the same with \(G_j(y)\) in place of \(G_j(y^-)\). This gives
\[
t_j=t_0+pF(p)[1-G_j(p^-)]+\int_p^\infty[1-G_j(y)][1-F(y)]\,dy
\]
and \(g_j\) in \eqref{eq:spread}. The difference of the two classes gives \(\DT\) in \eqref{eq:spread}. Each
integral is finite, since \(0\le\int_p^\infty[1-G_j(y)]\,dy\le\E_j\theta<\infty\) and \(0\le F\le1\).

Since \(\E_j[\max\{\theta,p\}]=p+\int_p^\infty[1-G_j(y)]\,dy\), the sum of the two lines of \eqref{eq:spread}
gives the split in its general form,
\[
\DT(F)+g_H(F)-g_L(F)=pF(p)D(p^-)+\E_H[\max\{\theta,p\}]-\E_L[\max\{\theta,p\}].
\]
The reserve term vanishes under equal admissibility, \(G_H(p^-)=G_L(p^-)\). It also vanishes when the incumbent
always meets the reserve, \(F(p)=0\).

\emph{Part (a).} In \eqref{eq:spread}, the weight \(1-G_j\) on \(F\) is nonnegative. The pointwise decrease
from \(F\) to \(\tilde F\) therefore lowers \(g_j\) by \(\int_p^\infty[1-G_j(y)][F(y)-\tilde F(y)]\,dy\ge0\).
The decrease is strict exactly when this integral is positive. Gross profit at a fixed belief \(\mu\) is
\((1-\mu)g_L+\mu g_H\), an average with fixed weights, so it falls as well.

\emph{Part (b).} The difference of \eqref{eq:spread} at the two distributions is
\[
\DT(\tilde F)-\DT(F)=\int_p^\infty[F(y)-\tilde F(y)]\,D(y)\,dy-p[F(p)-\tilde F(p)]\,D(p^-).
\]
Under the class ranking~\eqref{eq:classes}, \(D(p^-)=0\), so the second term is zero. Also \(D(y)\ge0\) and
\(F(y)-\tilde F(y)\ge0\) for every \(y\), so the integrand is nonnegative. The spread therefore rises, strictly
exactly when the integral is positive.

Equal admissibility also makes the second term zero. Together with \(D(y)\ge0\) for \(y>p\), it gives the same
conclusion. \(\square\)

\begin{named}{Proposition~1(c) (the weakest restriction)}
Let \(G_H\) and \(G_L\) have finite means. The spread weakly increases under every first-order stochastic
strengthening of every continuous \(F\) if and only if \(D(y)\ge0\) for every \(y>p\) and \(D(p^-)\le0\).
When neither class has an atom at \(p\), this condition is exactly first-order dominance above the reserve
together with equal admissibility, \(G_H(p^-)=G_L(p^-)\).
\end{named}

\paragraph{Proof of Proposition~1(c).}
Fix the incumbent's value \(R\). Suppose \(R\ge p\). The target receives \(\min\{R,\max\{\theta,p\}\}\) for every value of the rival. An admissible rival and the incumbent both bid at least \(p\), and an inadmissible
rival leaves the incumbent to pay \(p\). Suppose \(R<p\). The target receives \(p\1\{\theta\ge p\}\).

For \(R\ge p\), \(\E_j[\min\{R,\max\{\theta,p\}\}]=p+\int_p^R[1-G_j(y)]\,dy\). The expected difference in
proceeds between the classes at incumbent value \(R\) is therefore
\[
\Gamma(R)=
\begin{cases}
pD(p^-), & R<p,\\[3pt]
\int_p^RD(y)\,dy, & R\ge p,
\end{cases}
\]
and \(\DT(F)=\E_F[\Gamma(R)]\).

The function \(\Gamma\) is bounded, since \(\int_p^\infty|D(y)|\,dy\le\E_H\theta+\E_L\theta<\infty\). It is
constant below \(p\), continuous on \([p,\infty)\), and has derivative \(D\) almost everywhere above \(p\). It
moves from \(pD(p^-)\) to \(\Gamma(p)=0\) at \(p\). Thus \(\Gamma\) is nondecreasing if and only if \(D\ge0\) almost
everywhere above \(p\) and \(D(p^-)\le0\).

If \(\Gamma\) is nondecreasing, every first-order stochastic strengthening raises \(\E_F[\Gamma(R)]\). If it is not,
continuity on each side of \(p\) gives values \(0<R_1<R_2\), both different from \(p\), with
\(\Gamma(R_1)>\Gamma(R_2)\). Let \(F\) be uniform on \([R_1-\delta,R_1+\delta]\), and let \(\tilde F\) be the same law
shifted to \(R_2\). Then \(\tilde F\le F\), and as \(\delta\to0\),
\(\DT(\tilde F)-\DT(F)\to\Gamma(R_2)-\Gamma(R_1)<0\).

Suppose the support must start at zero. Mix the same small uniform component on \([0,\bar r]\) into both laws.
Both mixtures are continuous, and the mixture scales the difference by a positive factor and keeps its sign.

Distribution functions are right-continuous. Thus \(D\ge0\) almost everywhere above \(p\) is the same as
\(D(y)\ge0\) for every \(y>p\), and it implies \(D(p)\ge0\). When neither class has an atom at \(p\),
\(D(p^-)=D(p)\ge0\). The condition then holds exactly when \(D(p^-)=0\), which with the ranking above the
reserve is equal admissibility. \(\square\)

\paragraph{Scope of the class ranking.}
The proofs in this subsection use the class ranking~\eqref{eq:classes} only through \(D(y)\ge0\) for \(y>p\) and
\(D(p^-)=0\). These two facts give \(g_H\ge g_L\) and \(\DT\ge0\) in \eqref{eq:spread}, the comparison in part
(b), and the split identity of Section~\ref{sec:profits}. The same results therefore hold under the ranking
above the reserve together with equal admissibility.

For the point masses with \(p<\ell\), \(D(y)=\1\{\ell\le y<h\}\) and \(D(p^-)=0\). Also \(\Gamma(R)=(R-\ell)_+\)
for \(R<h\). Equation \eqref{eq:spread} then gives \(\DT=\int_\ell^h[1-F(y)]\,dy\), which equals
\(\int_\ell^{\bar r}[1-F(y)]\,dy\) at \(\bar r<h\). It also gives \(g_j=\int_p^\theta F(y)\,dy\) at \(\theta=h\)
and \(\theta=\ell\).

For the uniform law on \([0,r]\), these integrals reproduce \eqref{eq:a-payoffs}. The kernel \(\Gamma\) belongs to
the cash second-price rule with reserve. Section~\ref{sec:payment} changes the rule and the kernel with it.
